% !TEX program = pdflatex
\documentclass[12pt,a4paper]{article}

\usepackage[T1]{fontenc}
\usepackage[margin=1in]{geometry}
\usepackage{xcolor}
\usepackage{hyperref}
\usepackage{setspace}
\usepackage{parskip}
\usepackage{enumitem}

\hypersetup{
  colorlinks=true,
  urlcolor=blue!70!black
}

\onehalfspacing
\pagestyle{empty}

\begin{document}

\begin{center}
{\Large\bfseries Personal Statement}\\[4pt]
{\large Odeyemi Olusegun Israel}\\[2pt]
{\normalsize UK Global Talent Visa --- Digital Technology (Exceptional Promise)}
\end{center}

\bigskip

I am an independent researcher and systems engineer based in Derby,
United Kingdom.  Over the past three years---without institutional
affiliation, funding, or supervision---I have built a blockchain-native
compliance platform processing 2.64~million transactions, produced three
original research papers establishing a unified mathematical theory of
instability detection across finance, energy infrastructure, and fluid
dynamics, filed a UK patent application covering four inventions, and
open-sourced developer tools in TypeScript and Python.  I am applying
under the Exceptional Promise route because my work demonstrates the
capacity to become a leading figure in the UK's applied mathematics,
financial technology, and machine learning sectors.

%% ── SECTION 1: ENGINEERING ─────────────────────────────────────────

\textbf{1. What I have built.}

My central engineering project is \textbf{AMTTP}---the Anti-Money
Laundering Transaction Trust Protocol.  It addresses a problem that no
existing commercial product has solved: enforcing regulatory compliance
\textit{before} a blockchain transaction reaches finality, rather than
generating alerts after the fact.  Chainalysis, Elliptic, and TRM~Labs
all operate post-settlement; AMTTP intervenes at the protocol level,
within the approximately 12-second Ethereum block confirmation window.

The system spans four architectural layers: client SDKs (TypeScript and
Python, both open-sourced on GitHub), a compliance orchestration engine
combining ML risk scoring, graph-based relationship analysis, sanctions
screening, and configurable policy rules, an offline ML training pipeline,
and a smart contract layer with 18~contracts deployed on Ethereum Sepolia.
The infrastructure comprises 17~containerised microservices coordinated
through a central gateway.  The ML pipeline uses a Composite Teacher
architecture trained on 2.64~million transactions, achieving a production
ROC-AUC of \textbf{0.9879} and F1 of \textbf{0.9012} with 21~features and
zero training-serving skew.

I also built \textbf{zkNAF}, a zero-knowledge proof framework that allows
institutions to verify KYC credentials, risk score ranges, and sanctions
non-membership without exposing personal data---resolving the tension
between GDPR's data minimisation requirements and AML transparency
obligations at the architectural level.

The AMTTP architecture is described in my published paper:
\textit{Olusegun Odeyemi, ``AMTTP: A Four-Layer Architecture for
Deterministic Compliance Enforcement in Institutional DeFi,''
TechRxiv, February~27, 2026},
\href{https://doi.org/10.36227/techrxiv.177220113.33816607/v1}{DOI:\,10.36227/techrxiv.177220113.33816607/v1}.

%% ── SECTION 2: RESEARCH ────────────────────────────────────────────

\textbf{2. What I have discovered.}

Alongside the engineering work, I have pursued a research programme that
began with a practical question---why do fraud detectors miss what they
are designed to find?---and evolved into a unified mathematical theory of
instability in complex systems, with results spanning machine learning,
systemic financial risk, electrical grid failure, algorithmic stablecoin
collapse, and the regularity of the 3D Navier-Stokes equations.

\textit{Paper~1: Blind Spot Decomposition Theory (BSDT).}\;
Building on the engineering experience of AMTTP (see published paper
above), my first research paper provides a mathematical framework for
characterising missed fraud.  I decompose detection failures into four near-orthogonal
components---Camouflage~($\delta_C$), Feature Gap~($\delta_G$), Activity
Anomaly~($\delta_A$), and Temporal Novelty~($\delta_T$)---and demonstrate
that these account for over 95\% of explainable variance in missed-fraud
indicators.  The correction operators recover recall by up to
\textbf{84.4~percentage points} without retraining or access to model
internals.  Validated across 4~model architectures, 7~datasets from
6~domains, and over 36~model--dataset combinations, the best variant
achieves mean~F1 of~0.787 with 88.1\% false positive reduction.
Target journal: \textbf{IEEE~TKDE}.

\textit{Paper~2: Universal Deviation Law (UDL).}\;
My second paper tackles anomaly detection from first principles.
I introduce a multi-spectrum tensor framework with 14~composable
operators spanning statistical, geometric, spectral, and
information-theoretic families.  A $14\times14$ correlation audit
confirms that only 1~pair out of~91 is redundant: each operator captures
a genuinely distinct perspective on deviation.  The best configuration
achieves mean AUC of~\textbf{0.972} and 91\% coverage across five
benchmarks---categorically outperforming Isolation Forest~(0.788),
LOF~(0.632), and the current state-of-the-art ECOD~(0.921)---with a
single hyperparameter set and no per-dataset tuning.  I prove five
formal theorems providing local distinguishability, global sensitivity,
and finite-sample guarantees.
Target venue: \textbf{NeurIPS~2026}.

\textit{Paper~3: Systemic Risk as Geometry (AF+NS merged).}\;
My third paper---and the centrepiece of the research programme---establishes
a geometric framework linking the BSDT deviation operators to the
Lyapunov energy of networked agent dynamics.  It proves three theorems
(equivalence of the detection functional and the physical potential;
a spectral phase transition at a critical manifold; and a closed-form
optimal damping policy $\gamma^*(E) = E/(E+\theta)$) and validates
them across four domains:

\begin{itemize}[nosep, leftmargin=2em]
  \item \textbf{Banking} (140~quarters of FDIC call-report data):
    6-quarter lead before the 2008 GFC using a threshold trained
    exclusively on pre-2006 data; AUROC~$= 0.805$ at institution level;
    a 24-institution global panel spanning US, UK, EU, Asia, and Africa
    achieves AUROC~$= 0.811$ with a 19-quarter GFC lead.

  \item \textbf{Algorithmic stablecoins} (Terra/Luna UST, May 2022):
    detection at hour~23, one hour after the initial depeg event,
    with simulation-to-reality correlation $r = 0.996$.  Adaptive
    friction limits the depeg to 0.28\% (versus $>$99\% collapse
    without intervention).

  \item \textbf{Energy infrastructure} (Texas ERCOT, February 2021):
    detection at hour~24, a full $+$72~hours before rolling blackouts
    begin.  Adaptive friction reduces peak capacity loss by 71\%.
    The per-agent decomposition reveals that the cascade origin shifted
    from frozen wind turbines (the public narrative) to gas supply chain
    disruption (the structural vulnerability)---a finding invisible to
    event-study analysis.

  \item \textbf{Fluid dynamics} (3D incompressible Navier-Stokes):
    reinterpreting $\gamma^*$ as a state-dependent viscosity yields the
    \textbf{P/D halving principle}---adaptive viscosity halves the enstrophy
    production-to-dissipation ratio at every Reynolds number, confirmed
    numerically across $\mathrm{Re} \in [314, 6283]$.  This provides a
    new conditional regularity criterion connected to the Clay Millennium
    Prize Problem and identifies three strategies toward closing the gap
    to unconditional regularity.
\end{itemize}

\noindent
In all three non-financial domains, the dominant BSDT channel is
$\delta_C$ (camouflage)---the mechanism by which an unstable system
conceals its proximity to collapse---demonstrating that the mathematical
structure is genuinely universal, not fitted to any single domain.
Target journal: \textbf{SIAM Journal on Applied Mathematics}.

%% ── SECTION 3: IP ──────────────────────────────────────────────────

\textbf{3. Intellectual property.}

On 4~March~2026, I filed a UK patent application with the Intellectual
Property Office (Patents Act~1977) covering four inventions unified by
the blind-spot energy functional: AMTTP's pre-settlement compliance
architecture, the Adaptive Friction stabilisation mechanism, the
Universal Deviation Law's multi-spectrum tensor framework, and the BSDT
failure-mode correction operators.  The specification describes the
combined system as a \textit{``State-Adaptive Compliance Enforcement
System with Multi-Spectrum Anomaly Detection and Failure-Mode Correction
for Decentralised Finance Networks.''}

%% ── SECTION 4: IMPACT ──────────────────────────────────────────────

\textbf{4. Open-source contributions and tools.}

The full codebase---research code, pipeline implementations, client SDKs,
smart contracts, and replication packages for all three papers---is
publicly available on GitHub at
\texttt{github.com/segetii/fintech}.  Specific open-source components
include:

\begin{itemize}[nosep, leftmargin=2em]
  \item TypeScript and Python client SDKs for AMTTP integration;
  \item A MetaMask Snap for browser-based compliance verification;
  \item A GravityEngine simulation toolkit for networked agent dynamics;
  \item A 3D pseudo-spectral Navier-Stokes solver with BSDT diagnostic
        suite and adaptive viscosity;
  \item Complete replication packages for all empirical results.
\end{itemize}

%% ── SECTION 5: UK FIT ──────────────────────────────────────────────

\textbf{5. Why the UK, and why now.}

The UK is positioning itself as the global leader in responsible fintech
regulation.  The FCA's evolving framework for crypto-assets will create
substantial demand for compliance infrastructure that operates at the
protocol level.  AMTTP is, as far as I am aware from surveying the
commercial landscape, the only system that provides pre-settlement AML
enforcement on blockchain transactions.

The cross-domain reach of my research---from banking regulation to energy
grid resilience to pure mathematics---aligns with the UK's strategic
priorities in financial stability (Bank of England, PRA), infrastructure
resilience (Ofgem, National Grid ESO), and mathematical sciences (EPSRC,
the Heilbronn Institute).  The Navier-Stokes extension, in particular,
connects directly to the UK's world-leading position in applied
mathematics and fluid dynamics research.

My immediate plans are: (1)~submit the three papers to their target
journals (IEEE~TKDE, NeurIPS~2026, SIAM~J.~Applied Math);
(2)~prepare a short letter for \textit{Nature Communications}
summarising the cross-domain universality result for broader scientific
visibility; (3)~commercialise AMTTP as a compliance-as-a-service product
for UK-regulated financial institutions and virtual asset service
providers; and (4)~build a research-led company at the intersection of
machine learning, financial regulation, and mathematical stability
theory---headquartered in the UK.

%% ── CLOSING ────────────────────────────────────────────────────────

\medskip

I have done this work alone, without a university laboratory, a GPU
cluster, or a research budget.  In three years, I have gone from a
practical question about missed fraud to a unified mathematical framework
that detects collapse across banking systems, stablecoins, power grids,
and turbulent fluids---and I have built the production engineering to
deploy it.  What I bring is the ability to move between rigorous
mathematical theory and production-grade systems, and the discipline to
validate claims before making them.  I believe I can contribute
meaningfully to the UK's digital economy, scientific reputation, and
regulatory leadership, and I am asking for the opportunity to do so.

\end{document}
